\documentclass[11pt]{article}
\usepackage[T1]{fontenc}
\usepackage{amsmath}
\usepackage{amssymb}
\usepackage{array}
\usepackage{booktabs}
\usepackage[margin=1in]{geometry}
\usepackage[hidelinks]{hyperref}
\usepackage{seqsplit}
\setlength{\emergencystretch}{2em}

\newcommand{\claimstatus}[2]{\textbf{#1 [#2]}}
\newcommand{\poly}{\operatorname{poly}}
\newcommand{\bits}{\operatorname{bitlength}}
\newcommand{\archivepaper}{\path{paper/main.tex}}
\newcolumntype{L}[1]{>{\raggedright\arraybackslash}p{#1}}
\newenvironment{claimproof}{\par\noindent\textit{Proof.}}{\hfill\(\square\)\par}


\title{Cyclotomic Exceptional Extraction in\\
Signed Geometric-Sum Circuits}
\author{MOSEF Research Program}
\date{31 July 2026}

\begin{document}
\maketitle

\begin{abstract}
We classify the rational root-of-unity ratios arising from one signed
depth-two geometric-sum circuit and isolate two exceptional quadratic
cyclotomic families.  Beyond the common-step case, only order four with
ratio one and order six with ratio two survive.  Their quotient cofactors
have division-free logarithmic-step modular evaluators and can expose proper
factors after the direct cyclotomic, stage, and public-overlap GCDs are all
units.  Exact resultant formulas make every overlap branch public and
branch-total.  We also prove the quantifier limits: every fixed finite public
schedule misses infinitely many semiprimes, materializing exact lifts can
cost exponentially more than compact evaluation, and one cofactor supplies
only one signature bit.  These results concern a restricted circuit grammar;
they neither form a universal schedule nor solve general integer factoring.
\end{abstract}

\section{Charged circuit grammar}

Let
\[
 S_M(X)=1+X+\cdots+X^{M-1},\qquad
 Q_1(X)=S_A(X),\quad Q_2(X)=S_B(X^A),
\]
where public \(A,B\ge2\) are unequal.  For nonzero public coefficients set
\[
 F(X)=c_1Q_1(X)+c_2Q_2(X).
\]
The cost model charges construction, both stage evaluations, coefficient
reduction, every GCD, every failed inversion, independently evaluated
cyclotomic and cofactor exits, and every requested dense or sparse output.
Complexity is measured in
\(\bits(N)=\lfloor\log_2N\rfloor+1\).

\claimstatus{BAR-018}{PROVED}.  The stages are coprime over
\(\mathbb Q[X]\), with
\[
 \operatorname{Res}(Q_1,Q_2)=B^{A-1}.
\]
If \(L=Q_1(g)\) is a unit modulo \(N\), then
\[
 \gcd(F(g),N)=
 \gcd(c_1+c_2Q_2(g)L^{-1},N).
\]
A proper \(L\)-GCD already factors \(N\), while a full \(L\)-collision
reduces the aggregate to the public value \(c_2B\).  Moreover, for
\(h=\gcd(A-1,B-1)\),
\[
 Q_2(X)-Q_1(X)=XS_h(X)C_{A,B}(X),
\]
and its polynomial GCD with either natural endpoint is exactly \(S_h(X)\).

\begin{claimproof}
The identity
\[
 Q_2(X)-B=(X-1)Q_1(X)\sum_{j=1}^{B-1}S_j(X^A)
\]
gives coprimality and the resultant by evaluation at the \(A\)-th roots
other than one.  It also gives the three modular branches after evaluating
at \(g\).  Root-set comparison with \(X^{A-1}-1\) and
\(X^{B-1}-1\), followed by monic division, yields the exact common factor
\(S_h\).  Each modular denominator branch is therefore explicit; no silent
division is used.
\end{claimproof}

\section{Primitive residues and public overlap}

Write \(d=\gcd(|c_1|,|c_2|)\) and \(c_i=da_i\), with
\(\gcd(|a_1|,|a_2|)=1\).

\claimstatus{BAR-019}{PROVED}.  The content \(d\) has a total unit, proper,
or full GCD exit.  On the unit branch, the primitive aggregate
\(P=a_1Q_1+a_2Q_2\) has exact stage resultants
\[
 \operatorname{Res}(Q_1,F)=(c_2B)^{A-1},\qquad
 \operatorname{Res}(Q_2,F)=c_1^{A(B-1)}B^{A-1},
\]
and therefore
\[
 \gcd(Q_1(g),F(g),N)\mid\gcd(c_2B,N),\quad
 \gcd(Q_2(g),F(g),N)\mid\gcd(c_1B,N).
\]
At a root of unity away from both stages,
\[
 F(\zeta)=0\iff
 c_1(\zeta^A-1)^2+
 c_2(\zeta-1)(\zeta^{AB}-1)=0.
\]

\begin{claimproof}
Content factorization proves the first trichotomy.  Resultant
multiplicativity and evaluation on each stage root give the displayed
powers.  Bézout identities specialize modulo every prime divisor of \(N\),
so the overlap divides the public coefficient--multiplier GCD.  Multiplying
the rational stage equation by
\((\zeta-1)(\zeta^A-1)\) gives the final criterion.
\end{claimproof}

The boundary factors are not exhaustive.  With
\((A,B,c_1,c_2)=(3,7,1,1)\), \(\Phi_4\mid F\); at \(N=55,g=2\), both
stages and both public overlap bounds are units but
\(\gcd(F(2),55)=5\).

\section{Complete rational root-orbit classification}

For primitive \(\zeta\) of order \(n>1\), outside both stage zero sets, put
\[
 R_{A,B}(\zeta)=-\frac{S_B(\zeta^A)}{S_A(\zeta)}.
\]

\claimstatus{THM-003}{PROVED}.  The ratio is rational exactly in the
following cases:
\[
\begin{array}{c|c|c}
\text{order condition}&R_{A,B}(\zeta)&(c_1,c_2)\\ \hline
n\mid\gcd(A-1,B-1)&-1&(-1,1)\\
n=4,\ A\equiv B\equiv3\pmod4&1&(1,1)\\
n=6,\ A\equiv5,\ B\equiv3\pmod6&2&(2,1).
\end{array}
\]
Requested-order recognition is polynomial in the binary lengths of
\(A,B,n\).  The common-step family stays compactly described by the divisors
of \(\gcd(A-1,B-1)\); enumerating those divisors is separately charged.

\begin{claimproof}
Galois invariance and conjugation first force
\(n\mid A(B-2)+1\).  Choosing \(bA\equiv1\pmod n\) and
\(k\equiv b-1\pmod n\) gives
\[
 R_{A,B}(\zeta)+1
 =\left|\frac{1-\zeta^k}{1-\zeta}\right|^2
 =\frac{\sin^2(k\pi/n)}{\sin^2(\pi/n)}.
\]
This is a rational algebraic integer, hence a nonnegative integer.
Cyclotomic norms reduce the possibilities to the common-step orbit and the
two quadratic orbits.  Direct substitution gives the displayed ratios and
congruences.  The phase congruence alone is insufficient:
\((A,B,n)=(2,4,5)\) yields \((1+\sqrt5)/2\).
The complete norm case split is in
\path{research/proofs/THM-003-rational-root-orbits.md}.
\end{claimproof}

\section{Division-free exceptional cofactors}

Define
\[
\begin{aligned}
F_4&=S_A(X)+S_B(X^A)=\Phi_4(X)C_4(X),
 &&A\equiv B\equiv3\pmod4,\\
F_6&=2S_A(X)+S_B(X^A)=\Phi_6(X)C_6(X),
 &&A\equiv5,\ B\equiv3\pmod6.
\end{aligned}
\]

\claimstatus{BAR-020}{PROVED}.  Each \(C_i\) is a monic integer polynomial
of degree \(A(B-1)-2\), yet \(C_i(g)\bmod N\) has a constant-size descriptor
and can be evaluated without modular division in
\(O(\log A+\log B)\) modular operations.

\begin{claimproof}
For \(A=4a+3\), \(B=4b+3\), set
\[
 U_4(4t+3;x)=(1+x)S_t(x^4)+x^{4t},\quad
 D_4(A;x)=S_A(-x^2),
\]
and \(r=A-2\), \(s=A(B-2)\), \(k=(s-r)/2\).  Exact block division gives
\[
 C_4=U_4(A;x)+D_4(A;x)U_4(B;x^A)+x^rS_k(-x^2).
\]
For the order-six case, division by \(x^2-x+1\) has a period-six impulse
response; the quotient \(D_6(A;x)=\Phi_6(x^A)/\Phi_6(x)\) therefore has two
period-six segments.  With
\(H=x^3+2x^2+2x+1\), \(J=x^4+x^3-x-1\),
\[
\begin{split}
C_6={}&2H(x)S_{a+1}(x^6)\\
&+D_6(A;x)(H(x^A)S_b(x^{6A})+x^{6Ab})\\
&+2x^AJ(x)S_{Ab}(x^6).
\end{split}
\]
Binary geometric-sum evaluation proves the cost.  The integer identities
give additive prime-power valuations, so direct and cofactor GCDs remain
independently valid even on full-collision branches.  Dense output still
costs \(A(B-1)-1\) coefficients.
\end{claimproof}

\section{Overlap and schedule barriers}

\claimstatus{BAR-021}{PROVED}.  The cofactor GCD is proper exactly when its
capped local valuation profile is nonzero and not full.  The exact
stage/cofactor resultants have prime support in \(B\) for order four and
\(2B\) for order six.  Direct/cofactor overlap is controlled by
\[
\begin{aligned}
(u_4,v_4)&=\left(\frac{A(B+2)+1}{4},
                      \frac{A(B-2)+1}{4}\right),&
R_4&=u_4^2+v_4^2,\\
(u_6,v_6)&=\left(-\frac{2(A(B-2)+1)}3,
                      \frac{A(B+4)+4}3\right),&
R_6&=u_6^2+u_6v_6+v_6^2.
\end{aligned}
\]
Every fixed finite public joint schedule misses infinitely many square-free
semiprimes.

\begin{claimproof}
The valuation criterion follows from unique factorization.  Reducing the
cofactors modulo the removed quadratic gives \(u_i+v_iX\); the quadratic
norms are exactly \(R_i>0\), and resultant multiplicativity gives the stage
supports.  For a fixed schedule, multiply all positive charged exact values.
Only finitely many primes divide the product.  Any two distinct primes
outside that support form a semiprime on which every charged GCD is one.
The quantifier order fixes the schedule before the input.
\end{claimproof}

\claimstatus{BAR-022}{PROVED}.  For the balanced population
\[
 \mathcal P_m=\{p\text{ prime}:2^{m-1}\le p^2<2^m\},
\]
suppose materialized exact lifts have total bit budget \(W_m\) and hit
\(h_m\) of \(s_m=|\mathcal P_m|\) primes.  Exactly
\(\binom{s_m-h_m}{2}\) pairs make every materialized GCD one, and
\[
 \left\lfloor\frac{m-1}{2}\right\rfloor h_m\le W_m.
\]
Universal population coverage requires \(h_m\ge s_m-1\).

\begin{claimproof}
Pairs of unhit primes divide no lift and force every GCD to one.  The
square-free product of all hit primes divides the product of the nonzero
lifts, and every population prime has at least
\(\lfloor(m-1)/2\rfloor\) bits.  Taking binary lengths proves the budget.
This materialization theorem does not transfer to compact evaluation.
\end{claimproof}

\claimstatus{BAR-023}{PROVED}.  For the compact order-four tuple
\(A=3\), \(B_m=2^m+3\), \(g=2\), let \(C_m=C_4(2)\) and
\(E_m=3\cdot2^m+5\).  Then
\[
 C_m=\frac{16(2^{E_m}+3)}{35},\qquad
 \gcd(C_m,C_{m+1})=16.
\]
For a population of \(s\) odd primes, if exactly \(h\) divide \(C_m\), the
single GCD has \(h(s-h)\) proper outcomes,
\(\binom h2\) full collisions, and \(\binom{s-h}{2}\) unit outcomes.  It
cannot cover every pair when \(s\ge3\).

\begin{claimproof}
Substitution into the exact \(C_4\) identity gives the closed form.
Euclidean reduction of the two consecutive numerators gives the GCD 16,
after checking denominator primes 5 and 7 separately.  A pair splits exactly
when its two support bits differ, which gives the three counts and the
strict inequality \(h(s-h)<\binom s2\) for \(s\ge3\).
\end{claimproof}

\section{Reproduction and limitations}

\begin{verbatim}
python scripts/run_m24_rational_residue_audit.py
python scripts/check_m24_rational_residue_audit_differential.py
python scripts/run_m25_rational_root_orbit_audit.py
python scripts/check_m25_rational_root_orbit_differential.py
python scripts/run_m26_exceptional_cyclotomic_audit.py
python scripts/check_m26_exceptional_cyclotomic_differential.py
\end{verbatim}
The corresponding records are EXP-0023 through EXP-0025 under
\path{research/experiments}.  They validate finite symbolic and modular
instances; the proofs above are algebraic.

All full proofs, counterexamples, and charged branches remain in
\archivepaper.  Claim statuses and the focused projection are recorded in
\path{research/CLAIMS.md} and
\path{schemas/m82-paper-portfolio-v1.json}.  The paper proves results only
for the displayed signed geometric-sum grammar.  It supplies no universal
selector, no prime-support density theorem, and no general factoring
algorithm.

\end{document}
